\documentclass[11pt,a4paper]{article}
\usepackage{fontspec,polyglossia}
\setdefaultlanguage{ukrainian}
\setmainfont{DejaVu Sans}[ItalicFont=DejaVu Sans,ItalicFeatures={FakeSlant=0.2}]
\newfontfamily\cyrillicfont{DejaVu Sans}[ItalicFont=DejaVu Sans,ItalicFeatures={FakeSlant=0.2}]
\setmonofont{DejaVu Sans Mono}
\newfontfamily\cyrillicfonttt{DejaVu Sans Mono}
\usepackage{unicode-math}
\setmathfont{Latin Modern Math}
\usepackage[margin=22mm]{geometry}
\usepackage{amsmath,amsthm,hyperref}
\theoremstyle{definition}
\newtheorem{definition}{Означення}[section]
\newtheorem{proposition}[definition]{Твердження}
\setlength{\emergencystretch}{3em}
\sloppy\raggedright
\title{Лекція 5. Системи рівнянь у Python\\Математичні означення та критерії}
\author{Обчислювальна математика та програмування}
\date{}
\begin{document}
\maketitle

\section{Спільний розв’язок системи}
Нехай $D\subseteq\mathbb R^n$ та $f_i:D\to\mathbb R$, $i=1,\ldots,m$.
Система рівнянь має вигляд
\[
f_i(x)=0,\qquad i=1,\ldots,m,\qquad x\in D.
\]
\begin{definition}
Множина розв'язків системи є перетином множин розв'язків її рівнянь:
\[
Z=\{x\in D: f_1(x)=\cdots=f_m(x)=0\}
  =\bigcap_{i=1}^m\{x\in D:f_i(x)=0\}.
\]
Система сумісна, якщо $Z\ne\varnothing$, і несумісна, якщо $Z=\varnothing$.
Сумісна система визначена, якщо $Z$ містить рівно один елемент,
і невизначена, якщо має більше одного розв'язку.
\end{definition}
Дві системи еквівалентні на спільній області, якщо їхні множини розв'язків
збігаються. Перестановка рівнянь та множення окремого рівняння на
ненульову константу зберігають еквівалентність.

Для двох невідомих рівняння задають множини точок площини, а спільні
розв'язки є їхніми перетинами, включно з дотиками. Лінійна система
над $\mathbb R$ може мати нуль, один або нескінченно багато розв'язків.
Якщо $u\ne v$ є двома її розв'язками, то
\[
(1-t)u+tv,\qquad t\in\mathbb R,
\]
також є розв'язком; тому скінченної кількості більшої за один не буває.

\section{Матричний запис і єдиність}
Лінійна система з $m$ рівнянь та $n$ невідомих записується як
\[
Ax=b,\qquad A=(a_{ij})\in\mathbb R^{m\times n},\quad
x\in\mathbb R^n,\quad b\in\mathbb R^m,
\]
де рядок $i$ задає коефіцієнти $i$-го рівняння, стовпець $j$
відповідає невідомій $x_j$, а $b_i$ є правою частиною.
Відсутній у рівнянні доданок відповідає нульовому коефіцієнту.
Наприклад,
\[
\begin{cases}x+y=5,\\2x+y=7\end{cases}
\quad\Longleftrightarrow\quad
\begin{pmatrix}1&1\\2&1\end{pmatrix}
\begin{pmatrix}x\\y\end{pmatrix}
=\begin{pmatrix}5\\7\end{pmatrix}.
\]
\begin{proposition}
Для квадратної матриці $A\in\mathbb R^{n\times n}$ умови
\[
\det A\ne0,\qquad A\text{ оборотна},\qquad
Ax=b\text{ має єдиний розв'язок для кожного }b\in\mathbb R^n
\]
еквівалентні. Для фіксованого $b$ квадратна система має єдиний розв'язок
тоді й лише тоді, коли $\det A\ne0$.
\end{proposition}
Для порядку два
\[
\det\begin{pmatrix}a&b\\c&d\end{pmatrix}=ad-bc.
\]
У наведеному прикладі $\det A=1-2=-1\ne0$.
Умова $\det A=0$ сама по собі не розрізняє несумісність і нескінченність
множини розв'язків. Кількість рядків сама по собі не є критерієм сумісності.
Точне обчислення визначника для введених скінченних десяткових
коефіцієнтів можливе після представлення їх раціональними числами.

\section{Два способи розв’язання лінійної системи}
Нехай $A\in\mathbb R^{n\times n}$ і $\det A\ne0$.
Обернена матриця визначається рівностями
\[
A^{-1}A=AA^{-1}=I_n,
\]
де $I_n$ є одиничною матрицею. Помноживши $Ax=b$ зліва на $A^{-1}$,
одержуємо єдиний розв'язок
\[
x=A^{-1}b.
\]
Числове розв'язання $Ax=b$ може виконуватися без явного формування
$A^{-1}$. За точної арифметики обидва коректні способи дають той самий
розв'язок. У машинній арифметиці результати можуть відрізнятися через
округлення та властивості матриці.

Для кандидата $\widehat x$ вектор нев'язок і його максимальна норма:
\[
r=A\widehat x-b,\qquad
\|r\|_\infty=\max_{1\le i\le n}|r_i|.
\]
Навчальний критерій для зафіксованого масштабу даних:
\[
\|A\widehat x-b\|_\infty\le\varepsilon,
\qquad\varepsilon=10^{-6}.
\]
Кандидати двох способів можна зіставити за
$\|\widehat x^{(1)}-\widehat x^{(2)}\|_\infty$.
Зміна $b$ за сталого $A$ не змінює визначника й критерію єдиності,
хоча загалом змінює розв'язок. Для зміни $\Delta b$ маємо
$\Delta x=A^{-1}\Delta b$.

\section{Нелінійні системи та перевірка}
Нелінійна система записується у векторній формі
\[
F(x)=0,\qquad F:D\subseteq\mathbb R^n\to\mathbb R^n,
\qquad F=(f_1,\ldots,f_n)^{\mathsf T}.
\]
Для двох невідомих приклад має вигляд
\[
\begin{cases}x^2+y^2=5,\\y=x-1,\end{cases}
\qquad
F(x,y)=\begin{pmatrix}x^2+y^2-5\\y-x+1\end{pmatrix}.
\]
Його розв'язки $(2,1)$ і $(-1,-2)$ належать $D=\mathbb R^2$.
Якщо у формулі є знаменник $x^2+y^2$, точка $(0,0)$ виключається
з області. Усі інші обмеження складених виразів також зберігаються.

Початкове наближення для $n$ невідомих містить $n$ координат у
погодженому порядку. Локальний числовий пошук повертає одного
кандидата з даного старту; без додаткових припущень він не гарантує
збіжності, знаходження найближчого розв'язку чи повноти множини
знайдених розв'язків. Допустимість старту не гарантує допустимість
усіх проміжних ітерацій.

Для кандидата $\widehat z$ перевіряються скінченність дійсних координат,
належність $D$, доступна діагностика та
\[
\|F(\widehat z)\|_\infty
=\max_i|f_i(\widehat z)|\le10^{-6}.
\]
Для відповідних кандидатів двох інструментів навчальний критерій
узгодженості:
\[
\|\widehat z^{(1)}-\widehat z^{(2)}\|_\infty\le10^{-5}.
\]
Це критерії нев'язки й узгодженості, а не доведені межі похибки координат.
Для $F=Ax-b$ та точного розв'язку $x^*$ виконується
\[
\widehat x-x^*=A^{-1}r,
\qquad\|\widehat x-x^*\|_\infty
\le\|A^{-1}\|_\infty\|r\|_\infty,
\]
де $\|A^{-1}\|_\infty$ є максимальною сумою модулів елементів рядка.
Тому мала нев'язка без контролю властивостей матриці не дає такої самої
межі похибки невідомих.

Перевірка підстановкою у вихідні рівняння виявляє невідповідність
програмної моделі початковій постановці. Обчислення критеріїв виконується
до округлення для показу. Повторні кандидати того самого ізольованого
розв'язку не збільшують кількість різних розв'язків.

Для вікна $W\subseteq D$ знаходження коренів у $W$ саме по собі не
доводить відсутності коренів у $D\setminus W$. Для порівняння відстаней
до початку координат можна використовувати $\|z\|_2^2$;
для двох невідомих це $x^2+y^2$.
\end{document}
